\section{Introduction}

International shipping faces simultaneous pressure to reduce greenhouse-gas emissions, diversify its fuel base, and maintain operational reliability.
The IMO 2023 strategy and industry orderbook trends have turned alternative-fuel propulsion from a distant vision into a practical engineering reality
\cite{imo2023ghgstrategy,classnk2025alternativefuels}.
In this transition, fuels such as LNG, methanol, and ammonia are being introduced to lower carbon intensity
and to open potential opportunities toward net-zero shipping.
Consequently, marine dual-fuel and multi-fuel engines have become an important development direction for large propulsion systems
\cite{wingd_lng_faq_2025,wingd_ammonia_faq_2025,classnk_methanol_2024}.

However, alternative fuels introduce further challenges to marine-engine combustion processes.
To enable practical marine-engine operation with low-reactivity primary alternative fuels, an ignition source with high reactivity and reliability is still required;
without it, multiple important engine performance metrics, such as combustion stability, cycle-to-cycle repeatability, and load response can all deteriorate.
In other words, the fuel system may become lower-carbon and more flexible with fuel supply,
but the combustion process still depends on a small yet critical local event: pilot diesel injection,
often referred to as the pilot spray in marine engine context.

Pilot diesel injection therefore remains a critical enabling mechanism in multi-fuel marine engines:
it initiates combustion in a low-reactivity alternative-fuel main charge that does not auto-ignite under practical compression-ignition conditions
\cite{wingd_xdf_manual_2021,classnk_methanol_2024}.
In conventional diesel engines, diesel fuel is itself the main energy carrier
and compression ignition is already a mature engineering solution.
In multi-fuel engines, by contrast,
the pilot spray becomes a small but decisive ignition interface between the alternative main fuel and the combustion chamber.
For injector development, one central control problem is therefore the spatial development of the pilot spray itself:
spray penetration, cone angle, and the spatial morphology significantly influence the position and volume of the ignition kernel, and how it is positioned relative to the main-fuel charge.
That placement affects main-fuel mixing, flame propagation, and overall combustion efficiency
\cite{Karim2015DualFuel,WeiGeng2016DualFuelReview,Sahoo2009DualFuelReview}.
If pilot penetration is too short, the ignition kernel may remain too small or form in an unfavourable region with little main fuel;
if it is too long, the relative placement of the ignition kernel and main charge can also move away from the intended design, potentially leading to inefficient main-fuel burning and increasing the likelihood of wall-wetting as well as lube-oil dilution.
The penetration prediction pursued in this thesis is therefore not only a failure-avoidance tool, but a fast geometric descriptor for controlling pilot injection.

When that control objective is not met, the same pilot injection may become a source of operational risk instead.
If its liquid envelope reaches the piston crown or cylinder wall before a useful ignition kernel forms,
wall wetting, lube-oil dilution, degraded mixture preparation, and increased unburned-hydrocarbon emissions may follow
\cite{peraza2022swi,ahmad2025fueldilution,MorieiraMotaPanao2010}.
This risk is particularly relevant at the scale of large-bore marine engines,
where the nozzle-to-wall distance, spray momentum, and combustion-chamber dimensions together determine
whether liquid fuel can penetrate to near-wall regions and impinges on surfaces.
The practical question is therefore not only whether the pilot spray can ignite the main charge, and not only whether it hits a wall,
but whether its spatial development remains inside a controllable window that supports ignition, mixing, and safety margin.

%% Continue here
Existing research and engineering practice provide several established approaches for assessing
pilot-spray development and its relation to ignition placement and wall-interaction risk.
Engine tests are the closest to real operation, because they can examine spray, combustion, wall interaction,
and emissions within an integrated system.
High-fidelity CFD can resolve richer flow details and can be used to study coupling among spray formation, evaporation, mixing, and combustion.
Empirical penetration correlations and phenomenological models have also long been employed in diesel-spray research to describe spray development rapidly;
in this context, comparing penetration with wall distance provides a natural first-order geometric criterion for impingement screening
\cite{HiroyasuArai1990,NaberSiebers1996,Baumgarten2006MixFormation}.
Optical spray experiments occupy a complementary niche:
with a fixed optical setup, they can provide repeatable and directly comparable observations of the liquid envelope under various controlled conditions.
Among these techniques, high-speed Mie imaging is especially suitable for observing liquid-phase boundaries and macroscopic penetration evolution
\cite{Linne2013SprayImaging}.

However, each of these approaches covers only part of the need for fast, condition-parametric, experimentally grounded pilot-spray penetration prediction.
Engine tests are expensive, slow, and poorly suited to isolating spray development from a coupled engine system.
High-fidelity CFD is information-rich, but it requires complex inputs, model tuning, and substantial computational resources;
existing reacting diesel-spray LES studies show that a single higher-fidelity simulation can consume tens of thousands of CPU core-hours
\cite{CurtisNiemeyerSung2017,liu2020sprayreview}.
Empirical models are compact and interpretable, however, to cover multiple sources of complexity within the same numerical modelling scheme,
their functional forms often require additional empirical parameters, piecewise rules, or operating-condition-specific fitting,
which weakens rapid generalisation.
These complexitites typically include injector opening and closing transients, multihole plume-to-plume differences, and finite observation windows.
Optical experiments provide image-based evidence,
but the result of event-wise image post-processing is still essentially a tabulated dataset of discrete observations.
If left at that level, these observations cannot directly become a condition-parametric predictive model for fast design iteration.

Short pilot injections sharpen this tension.
Classical empirical penetration laws remain physically valuable,
but their applicability becomes more limited when the observable window is dominated by start-of-injection and end-of-injection transients
\cite{HiroyasuArai1990,ZhouLiYi2021,ZhouLiLaiWei2019}.
Spray-imaging literature also reminds us that penetration, cone angle, and boundary position are not method-independent absolute truths;
instead, they are operational observables produced jointly by the diagnostic method, thresholding rule, and post-processing definition,
meaning repeatable measurement-protocol-dependent quantities
rather than method-independent physical boundaries.
\cite{MagnottiGenzale2013,JohnsonNaberLee2012,ecn_liquid_penetration,ecn_liquid_penetration_accessory}.
The value of Mie imaging is therefore not that it directly provides a complete liquid-phase physical field,
but that, under clearly stated experimental and post-processing conditions,
it can provide stable geometric evidence relevant to engineering screening.

However, converting these geometric cues into surrogate training targets introduces several concrete difficulties.
A multihole nozzle produces several nominally similar plumes plumes but are not fully equivalent or uniform in direction, illumination, or penetration behaviour.
A fixed field of view turns part of the late-time penetration history into right-censored observations.
Raw Mie intensity is also affected by illumination, glare, and multiple scattering, while the extracted boundaries depend on thresholding and post-processing choices.
If the visible endpoint of each observed trajectory is treated directly as the true endpoint, the model can easily learn ``leaving the image''
as ``physical saturation''.
This means that thesis does not merely need more image data to improve the accuracy of a data-driven model,
nor does it need a more complicated end-to-end black-box regressor from images to geometric quantities.
The centraility of this thesis is an auditable workflow that connects imaging data, geometric definitions, censoring treatment, and uncertainty representation.

Data-driven surrogates offer a feasible direction for this goal.
Recent studies have shown that neural networks can predict macroscopic spray properties such as penetration and cone angle
from operating-condition parameters,
providing a predictive approach that is faster than high-fidelity CFD while more flexible than fixed-form empirical models
\cite{VazKarathanasis2026,liu2020sprayreview}.
In the present data setting, large-scale Mie imaging data provide an additional opportunity:
multi-campaign recordings acquired under the same optical configuration can be converted into consistently defined plume-resolved temporal observables
and used to train a mapping from conditions to penetration trajectories.

Nevertheless, the specific combination required for this workflow has received limited attention in the literature:
short-pulse pilot trajectories, right-censoring from a finite field of view, multihole plume registration, and uncertainty-aware outputs
that can be used directly in downstream geometric risk screening.
In other words, the missing piece is not a new theory of spray physics,
nor an additional isolated diagnostic measurement, but a fast, repeatable, explicitly scoped intermediate predictor step
between optical experiment archives and high-cost validation in engine-design workflow.

The aim of this thesis is therefore to develop and evaluate an auditable screening workflow
that converts a large archive of Mie-scattering spray videos into an early-stage engineering prediction layer
between optical experiments and high-cost validation based on engine tests.
To accomplish this aim, the thesis carries out four tasks:
\begin{enumerate}
  \item extract plume-resolved geometric observables from high-speed non-reacting Mie imaging videos;
  \item construct penetration-trajectory targets from finite-field-of-view observations
  \item train a surrogate that predicts both penetration mean together with conditional uncertainty; and
  \item connect the predicted distribution to a simplified representative engine geometry to form a probabilistic wall-impingement screening proxy.
\end{enumerate}

\subsection{Main contributions}

The main contributions of this thesis are as follows:
\begin{enumerate}
  \item a CUDA-accelerated workflow for converting terabyte-scale Mie-scattering videos into plume-resolved temporal observables and boundary trajectories;
  \item a reduced-order fitting strategy for trajectory alignment, robust extrapolation, and mitigation of right-censoring effects in penetration targets;
  \item a multi-stage uncertainty-aware learning framework that combines deterministic MSE training, heteroscedastic NLL training, and teacher-guided refinement on partially observed raw trajectories, culminating in an anchor-off Stage-3 surrogate with \SI{6.094}{mm} four-run mean RMSE under the external-cleaning protocol and conservative uncertainty coverage; and
  \item a simplified probabilistic impingement-screening method that decouples penetration prediction from downstream geometric collision evaluation while stating its prototype boundary explicitly.
\end{enumerate}

The scope of this work is bounded on purpose.
The thesis addresses non-reacting, non-evaporating pilot diesel spray under controlled pressurised-chamber conditions,
and the focus stays on the macroscopic liquid-envelope geometry and its temporal evolution as the optical setup sees it
--- not on a complete in-cylinder combustion process.

Several adjacent topics fall outside that scope by design: combustion chemistry, evaporation, dual-fuel ignition coupling,
thermal interaction between liquid spray and real engine surfaces, impingement on the piston-bowl inner wall,
and real-engine wall-film truth prediction. The model produced here should be read accordingly
--- as an early-stage screening surrogate, not as a finished truth model for engine wall impingement.

The MLP-based surrogate is also designed for interpolation within the operating-condition range that
the training data actually covers. For genuinely out-of-distribution extrapolation,
neither the predicted mean nor the uncertainty estimate is claimed at the same level of trust.
The model is therefore best understood as an engineering screening tool within its supported data domain,
not as a universally reliable predictor for arbitrary unseen conditions.

\subsection{Thesis structure}

The remainder of this thesis is organised as follows. Chapter~2 reviews the literature and establishes the technical rationale for the optical observables, target construction choices, and staged learning strategy. Chapter~3 describes the experimental data and image-processing workflow, covering the OSCC data, manual calibration, CUDA preprocessing, plume registration, strip extraction, CDF penetration, and boundary metrics. Chapter~4 describes trajectory targets, surrogate modelling, and probabilistic screening, covering trajectory cleaning, reduced-order fitting, feature design, Stage~1--3 learning, baseline construction, and the simplified wall-impingement proxy. Chapter~5 reports the empirical results, including the original held-out split, the external-cleaning protocol, Gaussian-process and physical baselines, CRPS and leave-one-nozzle-family diagnostics, campaign-level variation, and ablation evidence. Chapter~6 discusses why the final Stage-3 model is the most credible screening model, what has been validated, and where the current scope stops. Chapter~7 concludes the thesis by separating validated outcomes, current prototype boundaries, and the most important next steps.

